\section{What \odag{} must keep}
\label{sec:keep}

A shopping list is dangerous without a list of things that are not for sale.
This section names the seven commitments that give \odag{} its advantage, says
what each one is worth, and --- more usefully --- says what specific
temptation would destroy it. They are ordered by how expensive the mistake
would be.

\subsection{The cheap canonical form}
\label{sec:keep-canonical}

\paragraph{The commitment.} The stored form is the transitive reduction of the
asserted relation, and nothing else; values are canonicalised by denotation;
the serialisation is deterministic. Same knowledge, same bytes, same hash.

\paragraph{What it is worth.} Everything downstream. Content addressing,
snapshot identity, cheap diffing, structural sharing, agreement-by-hash,
convergent merge, and every certificate in Part~IV rests on this one property.
It is also unusual: neither a relational database nor a triple store nor a
graph database offers ``equal meaning, equal address''.

\paragraph{The temptation that would destroy it.} Expressiveness --- and the
forty-year record is worth putting beside it. The knowledge bases that
actually run the world are the \emph{weak} ones: SNOMED CT
\citep{spackman1997}, the Gene Ontology \citep{ashburner2000}, WordNet
\citep{miller1995}, schema.org \citep{guha2016} and Wikidata
\citep{vrandecic2014} are, in logical terms, little more than subsumption
plus vocabulary. Cyc \citep{lenat1995} bet on expressiveness plus hand-built
coverage; language models then commoditised the coverage, while the
expressiveness left the knowledge base hard to maintain and harder to share.
Every step
up the ladder --- arbitrary logical axioms, rules, defaults, weights ---
degrades the root from a fingerprint of \emph{knowledge} to a fingerprint of
\emph{phrasing}, because canonicalising up to logical equivalence means
canonicalising an entailment closure. The two-axis test is worth memorising:

\begin{keyidea}
A feature may enter the shared store only if it is
\textbf{(1) monotone} --- merge-as-union survives --- and
\textbf{(2) cheaply, semantically canonicalisable} --- ``equal knowledge,
equal root'' stays decidable and cheap.

Both are necessary. Neither is sufficient: features that pass both should
still wait for a real consumer.
\end{keyidea}

\subsection{Names as identity}
\label{sec:keep-names}

\paragraph{The commitment.} A node \emph{is} its name at every boundary. Names
are arbitrary and stable; meaning is carried by the structure around them, not
inside them.

\paragraph{What it is worth.} Merge is possible (\S\ref{sec:identity}),
references survive, and provenance has a subject. The contrast with \fca{},
where a concept is its extent--intent pair and therefore cannot survive an
edit, is the sharpest single difference between the two systems.

\paragraph{The temptation.} Meaning-bearing identifiers --- codes that encode
the classification, so that similar things get similar names. It is a very old
temptation and it has a very consistent record of failure
\citep{eco1995}:

\begin{center}\small
\begin{tabularx}{\textwidth}{@{}p{4.1cm} X@{}}
\toprule
\textbf{Attempt} & \textbf{How it failed}\\
\midrule
Wilkins's philosophical language (1668) \citep{wilkins1668}
  & A word's letters walk down the taxonomy --- so a concept with two
    legitimate genera has no legal name. \emph{The multi-parent problem,
    1668 edition.}\\
Dewey decimal (1876)
  & Meaning in the identifier bought real shelf locality, and paid for it with
    whole fields wedged into divisions chosen in the 1870s.\\
Linnaean binomials (1753)
  & One level of meaning in the name; taxonomic revisions therefore rename
    species, and synonymy databases exist to track the churn.\\
IUPAC names vs.\ CAS numbers
  & Chemistry has the purest meaning-bearing names anywhere --- and still
    keeps an arbitrary registry number beside them, because structures get
    re-perceived and databases need stable join keys.\\
\bottomrule
\end{tabularx}
\end{center}

There is also a hard mathematical reason, not merely a historical one: a
prefix or path code embeds a \emph{tree} order, and no linear key space can
embed a partial order with genuine multiple parents. A child would have to
share a prefix with both parents, which forces the parents to share it with
each other. So meaning-bearing names cannot represent a DAG at all --- which
is why Wilkins agonised over placements.

\paragraph{What to do instead.} Exactly what chemistry does: keep both. An
arbitrary stable name for identity, and a derived, regenerable, meaning-bearing
code (the ancestor-set code of Section~\ref{sec:dictionary}) as an index. That
settlement is available precisely because Proposition~\ref{prop:roundtrip}
says the code loses nothing.

\subsection{Assertion over derivation}

\paragraph{The commitment.} The store holds what was said. Consequences are
computed, never stored; nothing appears in the graph without an author.

\paragraph{What it is worth.} Attributability, bounded size, stability under
edit, and a coherent story about responsibility --- which becomes acute the
moment the writers are machines (\S\ref{sec:agents-review}).

\paragraph{The temptation.} Convenience. It is always tempting to materialise
a derived edge ``just this once'' for speed. The discipline that makes this
safe already exists and should be applied without exception: derived content
lives in a separate store with its own root, is never merged, and is
regenerable from the pinned root it was computed against. If derived content
must be shared, it enters as an ordinary claim by an identified author.

\subsection{One query primitive}

\paragraph{The commitment.} Conjunctive query is intersection of cones. There
is no second query mechanism.

\paragraph{What it is worth.} It is easy to underrate this, so here is the
concrete payoff. Because there is exactly one primitive:

\begin{itemize}[leftmargin=1.4em]
\item the planner is small enough to be \emph{proved} result-preserving, and
  is tested against a brute-force oracle over every one-, two- and three-term
  query on fixtures and random graphs;
\item typed values did not need a second mechanism --- a constraint like
  \verb|weight(..5kg)| is a query term, so matching is still one intersection,
  with one planner and one merge story;
\item disjunction is a thin layer that touches nothing stored;
\item consistency checking is the same primitive again: if disjointness
  claims are ever added, ``is anything both a Cat and a Dog?'' is
  \cmd{get(Cat, Dog)}. The check needs no checker.
\end{itemize}

\paragraph{The temptation.} A second query path for a special case ---
attribute payloads with a filter step, a text index consulted separately, a
graph traversal language. Each would double the surface that must agree with
itself.

\subsection{Mergeability}

\paragraph{The commitment.} Merge is union followed by re-reduction:
commutative, idempotent, convergent, always total.

\paragraph{What it is worth.} Multi-writer collaboration with no server, no
consensus, and no coordination --- and, as Section~\ref{sec:crypto} argues,
this is the property that makes \odag{} suit decentralised infrastructure
rather than merely tolerate it.

\paragraph{The temptation.} A merge that refuses. It is natural to want merge
to reject an inconsistent result --- two writers filing the same crate under
two different weights, say. But a merge that can fail is not total, and a
non-total merge is not a CRDT: replicas that reject different things diverge
permanently. The right answer is the one \odag{} settled on for constraints
generally: \emph{claims merge, enforcement is local}. The inconsistency
arrives, is visible, and is queryable --- and a reader may refuse to act on
it. Merge stays total.

\begin{pitfall}
An honest caveat, recorded because it is currently a documentation gap in the
system rather than a design decision: \odag{}'s \cmd{put} refuses to file one
item under two incompatible values of the same dimension, but \cmd{merge}
does \emph{not} enforce that guard --- so merging two stores can produce a
state \cmd{put} would have refused. Under the principle just stated this is
arguably correct (a refusing merge would not be total), but it should be
\emph{stated} rather than left implicit.
\end{pitfall}

\subsection{Exactness}

\paragraph{The commitment.} No floating point anywhere. Values are exact
rationals of an SI anchor unit; comparisons are exact arithmetic; there are no
weights, probabilities or confidences in the data.

\paragraph{What it is worth.} Determinism across replicas (G3), hence
canonical roots for typed values, hence the whole dimensions story. It also
has an unexpected downstream benefit: a structure with no floats and one query
primitive is unusually friendly to arithmetic circuits, which is what a
zero-knowledge proof would need (\S\ref{sec:crypto-zk}).

\paragraph{The temptation.} ``Just a confidence score.'' A number attached to
a claim seems harmless and is not: it breaks exact canonicalisation (which
precision? rounded how?) and breaks merge (whose number wins?). Confidence
belongs on a signed assertion record, where two authors disagreeing is
representable \emph{as two authors disagreeing} rather than as an averaged
scalar that belongs to nobody.

\subsection{One sort}

\paragraph{The commitment.} Everything is a node; ``category'' and ``item''
are roles, not types.

\paragraph{What it is worth.} Category hierarchies with no extra machinery,
and one traversal that crosses all levels (\S\ref{sec:one-sorted}).

\paragraph{The temptation.} Adding a type system --- classes versus
individuals, or a schema layer --- for the sake of validation. The cost is
that every such distinction becomes something two writers can disagree about,
and something the canonical form must encode.

\subsection{The meta-commitment: a place to put derived things}

Underneath the seven is one architectural habit that makes them affordable,
and it deserves a name because it is what allows Section~\ref{sec:learn}'s
shopping list to be so long.

\begin{keyidea}
\textbf{\odag{} has a well-defined place for derived structure}, with four
rules: it lives in a \emph{separate store with its own root}; it is
\emph{regenerable} from a pinned knowledge root; it is \emph{never merged};
and it is \emph{evictable} without data loss. Cone summaries, descendant
counts, ancestor-set codes, implication lints, MDL-learned concepts and
materialised meets are all instances.

This is what makes ``adopt, but not in the core'' a real answer rather than a
polite refusal. Because there is somewhere else to put things, \odag{} can be
generous about what it computes and strict about what it stores.
\end{keyidea}

On a storage network with rent, the pattern even acquires an economic
mechanism: a derived index is kept warm exactly as long as query traffic
justifies paying its postage, and is allowed to lapse otherwise
(\S\ref{sec:crypto-economics}). Regenerable things can be allowed to die.
